\documentclass[10pt,a4paper]{amsart}
\usepackage[margin=19mm]{geometry}
\usepackage{amsmath,amssymb,mathtools}
\usepackage[scr=rsfs,cal=euler]{mathalfa}
\usepackage{xfrac}
\usepackage[unicode,hidelinks,bookmarks=false]{hyperref}
\newcommand{\ie}{i.e.\ }
\newcommand{\cf}{cf.\ }
\newcommand{\resp}{respectively\ }
\newcommand{\Subsection}{\subsection}
\providecommand{\nobreakdash}{\nobreak\hbox{-}}
\setlength{\emergencystretch}{2em}
\multlinegap0pt
\usepackage{bm}
\usepackage{amssymb}
\usepackage{xcolor}
\usepackage{float}
\usepackage{tikz}
\usepackage[shortlabels]{enumitem}
\usepackage{mathtools}
\newtheorem{theorem}{Theorem}
\newcommand{\CSS}{\bar{X}}
\newcommand{\NN}{\mathbb{N}}
\def\be{\begin{equation}}
\def\del{\partial}
\def\ee{\end{equation}}
\def\lp{\left(}
\def\rp{\right)}
\title{Multiple diffusion scales and diffusion-driven instability: Emergence of near- and far-from-equilibrium patterns}
\author{Th\'eo Andr\'e, Szymon Cygan, Anna Marciniak-Czochra, Finn M\"unnich}
\date{Source: arXiv:2511.15648v2 (CC BY 4.0)}
\begin{document}
\maketitle

\noindent\emph{Source and licence.} Multiple diffusion scales and diffusion-driven instability: Emergence of near- and far-from-equilibrium patterns. Version arXiv:2511.15648v2 (CC BY 4.0), distributed by arXiv under the Creative Commons Attribution 4.0 International licence, http://creativecommons.org/licenses/by/4.0/, as recorded in the arXiv licence metadata for this exact version and shown on the abstract page https://arxiv.org/abs/2511.15648v2. Full licence notice: \texttt{licenses/arxiv-2511.15648-CC-BY-4.0.txt}.\par\smallskip

\noindent\textit{Editorial scope.} Counted unit: the article's theorem lem:DDI\_toysys\_exact (source0.tex lines 1271--1283). The article's complete original proof of it is delivered with it (source0.tex lines 1285--1307, one \texttt{proof} environment of 23 lines). No mathematics is written, repaired or re-derived: the statement and the proof are the article's own lines, in the article's own notation, and no retained line is substituted or re-flowed.  Retained uncounted as the source-local context the proof needs: lines 1155--1167.  Nothing was omitted from any retained span, so the package carries no omission record.  No retained line contains a citation, so the wrapper carries no bibliography and no bibliography entry.  No retained line contains a figure environment, an image-inclusion command or drawing code, so no image is delivered and the package declares no asset dependency.  Editorial formatting only, in addition: an AMS \texttt{amsart} preamble that restates the article's own declarations and macros used by the retained lines, namely the environments \texttt{theorem} on separate counters, together with the standard packages those lines need.  The retained environments print in this document as Theorem 1 in order of appearance, whereas the article prints the same items with the numbering of its own full document.  The complete native source of the article as distributed in the arXiv e-print for this exact version is bundled unchanged as \texttt{sources/189466/source0.tex}.
Let $\Omega = (0,L)\subset\mathbb{R}$, with $L>0$. The governing equations are
\be \label{eq:toysys}
\left\{
\begin{aligned}
    \del_t u &= -\mu_1 u + m_1 uv (1 + u v)^{-1}, & \textup{for } \ x \in \bar\Omega, \ t > 0,\\
    \del_t v &= D_v \Delta v -\mu_2 v + m_2 uv (1 + uv)^{-1} - vw, & \textup{for } \ x \in \Omega, \ t > 0,\\
    \del_t w &= D_w \Delta w -\mu_3 w + m_3 uv (1 + uv)^{-1}, & \textup{for } \ x \in \Omega, \ t > 0,
\end{aligned}
\right.
\ee
supplemented with homogeneous Neumann boundary conditions for $v$ and $w$, and nonnegative initial data.
Here $\mu_i \in (0,1]$ denotes decay rates, the $m_i \in [2,+\infty)$ are production rates, and $D_v, D_w > 0$ are the diffusion coefficients of ligands and enzymes respectively. The receptor population, $u$, is assumed diffusion-free.
The nonlinear production terms follow a Michaelis--Menten type saturation in the ligand--receptor binding.

\begin{theorem}
\label{lem:DDI_toysys_exact}
    Let $\Omega = (0, L)$ and consider the constant steady state $\CSS_+$ of~\eqref{eq:toysys}. This steady state exhibits DDI if either of the following conditions is satisfied:
    \begin{enumerate}
        \item[(I)] $L >0$, the diffusion coefficient $D_w > 0$ is fixed and 
        \begin{align} 
        0 < D_v < \varepsilon:=\sup_{j\in\NN}\dfrac{\det(J) - \det(J_{12}) \, D_w \lambda_j}{\lambda_j \lp \det(J_{13}) - J_1 \, D_w \lambda_j \rp}, 
        \end{align}
        \item[(II)] $D_v > 0$ fixed, and
        $$ L > \pi j \sqrt{\dfrac{J_1 \, D_v }{\det(J_{12})}}, \qquad D_w >  \dfrac{\det(J) - \det(J_{13})\,D_v \lambda_j}{\lambda_j(\det(J_{12}) - J_1\,  D_v \lambda_j)} > 0 $$
        holds for at least one $j \in \mathbb N$.
    \end{enumerate} 
\end{theorem}

\begin{proof}
    To determine the precise parameter regimes for DDI, we analyze the sign of the Routh--Hurwitz element $p_3(\lambda_j) = - \det(J - \lambda_j D)$. Instability requires $p_3(\lambda_j)<0$, and we distinguish the cases of small $D_v$ and large $D_w$ accordingly.
    
    By solving $p_3(\lambda_j) < 0$ for $D_v$, that is, finding values of $D_v$ such that
    \begin{align}
    -\det(J) + \biggl( \det(J_{13})\,D_v + \det(J_{12}) \,D_w \biggr) \lambda_j + \biggl( -J_1\, D_v D_w \biggr) \lambda_j^2 < 0,
    \end{align}
    we obtain condition~\emph{(I)}.  
    The right-hand side of the resulting inequality is positive for sufficiently large $j \in \mathbb{N}$ and converges to $0$, since $J_{12}$ is unstable while $J_1$ and $J_{13}$ are always stable.

    The second condition is obtained by fixing $j \in \mathbb{N}$ and rewriting $p_3(\lambda_j) < 0$ as
    \begin{align}
    \label{eq:toy_critical_d2}
        \lambda_j \big( \det(J_{12}) - J_1\, D_v \lambda_j \big) D_w 
        < \det(J) - \det(J_{13})\,D_v \lambda_j.
    \end{align}
    Since the right-hand side of~\eqref{eq:toy_critical_d2} is negative, this inequality can only hold for large $D_w$ if
    \begin{align}
        \det(J_{12}) - J_1 D_v \lambda_j < 0.
    \end{align}
    The latter can always be achieved by increasing the domain length $L$ if necessary. 
    Solving~\eqref{eq:toy_critical_d2} for $D_w$ and the above inequality for $L$ yields the two expressions stated in condition~\emph{(II)}.
\end{proof}

\end{document}
